\documentclass[11pt]{article}
\usepackage[T1]{fontenc}
\usepackage{lmodern,amsmath,amssymb,amsthm,mathtools,booktabs,microtype}
\usepackage[margin=1in]{geometry}
\usepackage[hidelinks]{hyperref}
\usepackage{xcolor}
\usepackage{fancyhdr}
\usepackage{url}
\definecolor{ink}{RGB}{24,47,72}
\definecolor{link}{RGB}{31,91,139}
\hypersetup{colorlinks=true,linkcolor=link,urlcolor=link,citecolor=link}
\pagestyle{fancy}\fancyhf{}\lhead{\small Quantitative GRH-method transfers}\rhead{\small Research note}\cfoot{\thepage}\setlength{\headheight}{14pt}
\newtheorem{theorem}{Theorem}[section]
\newtheorem{lemma}[theorem]{Lemma}
\newtheorem{corollary}[theorem]{Corollary}
\newtheorem{proposition}[theorem]{Proposition}
\theoremstyle{remark}\newtheorem{remark}[theorem]{Remark}
\newcommand{\Q}{\mathbf Q}\newcommand{\Z}{\mathbf Z}\newcommand{\C}{\mathbf C}\newcommand{\Ocal}{\mathcal O}\newcommand{\eps}{\varepsilon}\newcommand{\vp}{v_p}\newcommand{\rh}{r_h}\newcommand{\Gauss}{\mathsf G}\newcommand{\pbar}{\overline{\Q}_p}
\title{\color{ink}\bfseries Beyond Dirichlet GRH:\\ Polynomial Order-Exponent Bounds for\\ Simultaneous \(p\)-adic Twisting}
\author{Research analysis and proposed extension of \texttt{lsernau/grh-lean}}
\date{25 November 2026}
\begin{document}
\maketitle
\begin{abstract}
The Alp\"oge--Claude proof of RH and Dirichlet GRH, as rewritten in
\texttt{lsernau/grh-lean}, uses a simultaneous torsion-point selection
lemma to construct anticyclotomic twists with small Selmer lengths and
controlled special values. Its proof chooses an upper order exponent via
a divergent harmonic sum. We prove a different quantitative estimate:
for a fixed finite packet of integral polynomials, a sufficient upper
exponent is polynomial in inverse valuation tolerance and the valuation
budgets, independently of polynomial degrees. We also prove that a
positive fraction of the primitive roots in at least one admissible shell
can satisfy the entire packet. The mechanism is a Gauss-valuation saving
from roots on earlier shells, used in a multiplicative potential argument.
We transfer the estimate to multivariate polynomials, bounded power
series, and the repository's arithmetic twisting construction, obtaining
arbitrarily large simultaneous packets of characters of a common tame
conductor, with an explicit order-exponent bound. This is a mathematical
strengthening of the method, not a standard consequence of the final GRH
theorem. No priority claim or resolution of a named major open conjecture
is made; the arithmetic conclusions are restricted to the hypotheses of
the existing twisting theorem.
\end{abstract}
\tableofcontents
\newpage
\section{Source audit and research target}

The repository's top-level README describes two editions of the proof,
with a common set of five conclusions and a machine-checking history. At
upstream commit \texttt{72fead102f66} (full hash in the source audit), the
statements are written using Mathlib's Riemann zeta and Dirichlet
$L$-functions. In particular, the half-plane theorem quantifies over
Dirichlet characters of every positive modulus, not just primitive
characters; the principal pole is excluded. The readable edition has
approximately 300,000 lines and the compact edition about 76,000 lines
of code, as reported by their README files. The repository reports kernel
replay and checks with only \texttt{propext}, \texttt{Classical.choice}
and \texttt{Quot.sound}. We did not independently build or replay that
proof. The result taken as input here is a specific intermediate twisting
construction, not a claim that all automorphic, Artin or Dedekind
$L$-functions satisfy GRH.\footnote{Repository: \url{https://github.com/lsernau/grh-lean}. See
\texttt{README.md}, \texttt{grh-compact/README.md}, and
\texttt{grh/src/Statements.lean}.}

The compact proof guide and \texttt{grh-compact/src/New/Main.lean}
describe an arithmetic-to-analytic chain: short integral forms, a
spanning argument, a lower height bound for integral forms, and a
comparison of corrected covolumes. A common main height cancels, giving
an inequality for a prime remainder; an analytic one-sided criterion then
excludes off-line Dirichlet zeros. The Selmer twisting argument is part of
the lower height bound and is logically useful independently of the
prime-remainder conclusion.

\paragraph{The selected interface.}
The readable edition's \path{New/Congruences/Selmer/RootsOfUnity.lean}
constructs roots of unity simultaneously controlling polynomial
valuations. The general arithmetic theorem in
\path{New/Congruences/Selmer/GeneralTheorem.lean} controls Selmer lengths,
Stickelberger values, and local tame-generator valuations for the first
$m$ powers of one anticyclotomic character. The proof first bounds the
number of variables, then chooses an order exponent before the
polynomials are supplied. Keeping this quantifier order is essential:
choosing a conductor after seeing the polynomials is not a replacement
for the upstream selection step.

\paragraph{Relationship to earlier ProveIt reports.}
We listed the 21 entries in \path{docs/incoming} through GitHub and
inspected the LaTeX of two particularly relevant archives:
\path{grh_method_transfers_2026-11-25(1).zip}, sections
\path{torsion_transfer.tex} and \path{torsion_literature.tex}; and
\path{grh_method_transfers_2026-11-25.zip}, section
\path{ramification_budget.tex}. The former extends the harmonic-sum
selector to bounded analytic data, including infinite families subject
to a total budget. The latter identifies the growing ramification index
as an obstruction to rounding small normalized Selmer lengths to zero,
and extracts bounded support size from the arithmetic construction.
We also inspected the manuscript inventory of three other archives.
These are not claimed to be an exhaustive survey of every archive or
of the literature. The present contribution is a different conductor-order
estimate, with a density refinement for a fixed finite packet; it
complements rather than supersedes the cardinality-free results of the
prior report. See \texttt{SOURCE\_AUDIT.md} for pinned references.

\paragraph{Related open problems.}
Quantitative simultaneous central-value nonvanishing for finite-order
twists (including the higher-order twist problem of David--Fearnley--Kisilevsky),
control of Selmer groups in character families, and specialization of
Iwasawa-theoretic determinant ideals are natural targets. Horizontal
$p$-adic $L$-functions already give important simultaneous nonvanishing
results \cite{KN}. The theorem below supplies an improved quantitative
local ingredient for a finite packet of integral interpolation
polynomials or bounded power series. It does not supply an interpolation
identity or an initial nonzero witness for a specific arithmetic family.

\section{A degree-independent polynomial exponent bound}

Fix a prime $p$, and use the valuation $v$ on $\pbar$ normalized by
$v(p)=1$, with $v(0)=+\infty$. Let $\Ocal$ denote its ring of integers.
For $h\ge1$, define
\[
 e_h=\varphi(p^h)=(p-1)p^{h-1},\qquad r_h=e_h^{-1},\qquad
 \mathcal U_h=\{\zeta:\operatorname{ord}(\zeta)=p^h\}.
\]
Each element of $\mathcal U_h$ has $v(\zeta-1)=r_h$. Write
$c=1-1/p\in(0,1)$. For a nonzero polynomial $g(T)=\sum_j a_jT^j$,
let
\[
 \Gauss_\rho(g)=\min_j\{v(a_j)+j\rho\}\quad (\rho\ge0).
\]
For integral coefficients, $\Gauss_\rho(g)\ge0$. For $v(y)\ge\rho$,
$v(g(y))\ge\Gauss_\rho(g)$. Multiplicativity of the Gauss valuation
and factorization give, if $g=a\prod_{\alpha}(T-\alpha)$ with roots
counted with multiplicity,
\begin{equation}\label{eq:gauss}
 \Gauss_\rho(g)=v(a)+\sum_{\alpha}\min\{\rho,v(\alpha)\}.
\end{equation}
The formula remains valid for roots of negative valuation and roots
at zero. For completeness, multiplicativity follows by choosing, in
each factor, the smallest index attaining the minimum weighted
coefficient valuation. Their sum index in the product has a unique
term of minimal weighted valuation: any other pair with the same sum
index has an index below one of the chosen minima, or a strictly larger
weighted valuation. Thus that coefficient cannot cancel. Apply this to
the linear factors, for which the formula is immediate.

\begin{theorem}[Polynomial shell bound and simultaneous good fraction]\label{thm:one}
Let $g_1,\ldots,g_n\in\Ocal[T]$ be nonzero, $n\ge1$, let $\rho>0$,
and let $B_i\ge0$, $\delta_i>0$ satisfy
\[
 A_i:=\Gauss_\rho(g_i)\le B_i.
\]
Let $h_0\ge1$ satisfy $r_{h_0}\le\rho$, and let $0<\eta\le1$. Define
\begin{equation}\label{eq:constants}
 C=1+\max_i(B_i+\delta_i),\qquad
 P=\prod_{i=1}^n\left(1+\frac{cB_i}{\delta_i}\right).
\end{equation}
Choose an integer $H\ge h_0$ such that
\begin{equation}\label{eq:Hcondition}
 \left(\frac{H+C+1}{h_0+C}\right)^{c\eta}\ge P.
\end{equation}
Then there exists $h_0\le h\le H$ such that strictly more than a fraction
$1-\eta$ of $\zeta\in\mathcal U_h$ satisfy \emph{all} the inequalities
\begin{equation}\label{eq:desired}
 v(g_i(\zeta-1))\le A_i+\delta_i\qquad(1\le i\le n).
\end{equation}
In particular, all these values are nonzero. One explicit sufficient
choice is
\begin{equation}\label{eq:H}
 H=\left\lceil (h_0+C)P^{1/(c\eta)}\right\rceil.
\end{equation}
The bound does not depend on degrees, roots, or coefficient fields.
\end{theorem}

When $\eta=1$ the assertion is existence; when $\eta=1/2$ more than
half of the primitive roots in the selected shell work. It is important
that we compare with $\Gauss_\rho$, not with $\Gauss_{r_h}$. The latter
is a stronger requirement; the upstream harmonic-sum lemma establishes
that stronger comparison. Its strength is unnecessary for the subsequent
witness-based twisting theorem.

\subsection{Capped shell averaging}
\begin{lemma}[Nearest primitive root]\label{lem:nearest}
For $v(\alpha)=r_h$ and $C>0$, put
\[
 E_\alpha(\zeta)=v(\zeta-1-\alpha)-r_h\in[0,+\infty].
\]
Then
\[
 \frac1{e_h}\sum_{\zeta\in\mathcal U_h}\min\{E_\alpha(\zeta),C\}
 \le r_h(h+C).
\]
\end{lemma}
\begin{proof}
Choose $\zeta_*\in\mathcal U_h$ maximizing $v(\zeta-(1+\alpha))$.
For $\zeta\ne\zeta_*$, the ultrametric inequality gives
\[
 v(\zeta-(1+\alpha))\le v(\zeta-\zeta_*).
\]
Indeed the valuation of the difference of $\zeta-(1+\alpha)$ and
$\zeta_*-(1+\alpha)$ is at least the smaller valuation. The summand at
$\zeta_*$ is at most $C$; every other summand is at most
$v(\zeta-\zeta_*)$. Since
\[
 \prod_{\zeta\in\mathcal U_h\setminus\{\zeta_*\}}(\zeta_*-\zeta)
 =\Phi'_{p^h}(\zeta_*),
\]
we obtain a sum bounded by $C+v(\Phi'_{p^h}(\zeta_*))$. From
$\Phi_{p^h}(X)=(X^{p^h}-1)/(X^{p^{h-1}}-1)$, its derivative at a
primitive root has valuation $h-1/(p-1)\le h$.
The valuation of a primitive $p$-th root minus one is $1/(p-1)$, by
$\Phi_p(1)=p$ and equality of the valuations of all its factors.
The same argument with $\Phi_{p^h}(1)=p$ gives $r_h$ for the primitive
$p^h$-th roots. Dividing by $e_h$ completes the proof. The argument
includes $p=2,h=1$, where the derivative valuation is zero.
\end{proof}

\subsection{The saving from earlier shells}
For polynomial $g_i$, let $R_{ih}$ denote the multiset of roots of
valuation $r_h$, and define
\[
 w_{ih}=r_h\#R_{ih},\qquad
 W_{i,h}=\sum_{k=h_0}^h w_{ik},\qquad W_{i,h_0-1}=0.
\]
By subtracting the formula \eqref{eq:gauss} at radius zero from the
formula at $\rho$, we get
\begin{equation}\label{eq:total}
 W_{i,H}\le \Gauss_\rho(g_i)-\Gauss_0(g_i)\le A_i\le B_i.
\end{equation}
At shell $h$, define the slack $s_{ih}=A_i-\Gauss_{r_h}(g_i)$. All
contributions to this difference are nonnegative, even from nonintegral
roots. Every root on an earlier admissible shell $k<h$ contributes
$r_k-r_h\ge c r_k$. Consequently
\begin{equation}\label{eq:slack}
 0\le c W_{i,h-1}\le s_{ih}\le B_i.
\end{equation}
For $\zeta\in\mathcal U_h$, all roots not on shell $h$ contribute
exactly their Gauss-valuation term, since their valuations differ from
$v(\zeta-1)$. Thus, including the case of a zero evaluation,
\begin{equation}\label{eq:exact}
 v(g_i(\zeta-1))=\Gauss_{r_h}(g_i)+\sum_{\alpha\in R_{ih}}E_\alpha(\zeta).
\end{equation}
A root on an earlier shell has not simply used up part of the budget;
its contribution is now smaller by a definite proportion. This is the
extra information discarded by the harmonic-sum comparison.

\subsection{Multiplicative potential proof of Theorem \ref{thm:one}}
Let $D_{ih}(\zeta)=\sum_{\alpha\in R_{ih}}\min(E_\alpha(\zeta),C)$.
If inequality \eqref{eq:desired} fails for $i$, formula \eqref{eq:exact}
implies that the uncapped excess exceeds $s_{ih}+\delta_i$. Since
$C>B_i+\delta_i\ge s_{ih}+\delta_i$, the capped sum still exceeds it:
either one term reaches the cap, or no term changes. In particular,
\[
 \frac{D_{ih}(\zeta)}{\delta_i+cW_{i,h-1}}>1
 \quad\text{whenever constraint $i$ fails.}
\]
If a fraction at least $\eta$ of the shell fails at least one constraint,
then the average over the shell of the sum of these nonnegative ratios
is strictly greater than $\eta$. Lemma \ref{lem:nearest} bounds that
average from above, yielding
\begin{equation}\label{eq:bad}
 \sum_i\frac{w_{ih}}{\delta_i+cW_{i,h-1}}>
 \frac{\eta}{h+C}.
\end{equation}
The strict inequality is valid also at $\eta=1$ because every bad
point has a ratio sum strictly greater than one. Set
\[
 \mathcal P_h=\prod_i(1+cW_{i,h}/\delta_i).
\]
If shell $h$ has bad fraction at least $\eta$, then
\begin{align*}
 \frac{\mathcal P_h}{\mathcal P_{h-1}}
 &=\prod_i\left(1+\frac{cw_{ih}}{\delta_i+cW_{i,h-1}}\right)\\
 &\ge1+\sum_i\frac{cw_{ih}}{\delta_i+cW_{i,h-1}}
 >1+\frac{c\eta}{h+C}.
\end{align*}
Suppose all shells $h_0,\ldots,H$ have bad fraction at least $\eta$.
Since $\mathcal P_{h_0-1}=1$, and $0<c\eta\le1$, concavity of
$u\mapsto u^{c\eta}$ gives
\[
 \mathcal P_H>\prod_{h=h_0}^{H}\left(1+\frac{c\eta}{h+C}\right)
 \ge\left(\prod_{h=h_0}^H\frac{h+C+1}{h+C}\right)^{c\eta}
 =\left(\frac{H+C+1}{h_0+C}\right)^{c\eta}.
\]
But \eqref{eq:total} implies $\mathcal P_H\le P$, contradicting
\eqref{eq:Hcondition}. Thus there is a shell with bad fraction less
than $\eta$, as claimed. \qed

\subsection{Comparison with the harmonic estimate}
The upstream one-polynomial sufficient condition is
\[
 B<\delta\sum_{h=h_0}^H(h+1)^{-1},\qquad 0<\delta<1.
\]
Here we can take $B$ to be the total budget for the product of a finite
packet, using multiplicativity and then controlling each individual
loss relative to its own witness. Because the sum is at most
$\log((H+1)/h_0)$, this sufficient condition requires
$H+1>h_0\exp(B/\delta)$. It proves a comparison to the smaller
$\Gauss_{r_h}$; hence it is not an inferior proof of the exact same
stronger assertion. Our estimate addresses the weaker comparison which
is sufficient for the arithmetic application.

For equal budgets $B_i\le B_*$ and tolerances $\delta_i=\delta$, our
bound is of size
\[
 H=O\big((h_0+B_*+\delta+1)
              (1+cB_*/\delta)^{n/(c\eta)}\big).
\]
This is polynomial in $1/\delta$ for a fixed finite packet and fixed
$p,\eta,h_0,B_*$. The character order itself is bounded by $p^H$,
not by a polynomial. In the small tolerance regime this changes the
available estimate from a double-exponential order bound to a
single exponential of a polynomial. It is not a uniform improvement
when $n$ grows, so retaining both selection lemmas is prudent.

\begin{table}[ht]
\centering\small
\begin{tabular}{rrrrrrr}
\toprule
$p$ & $n$ & $B_*$ & $\delta$ & $\eta$ & $H$ new & $\log_{10}(H_{\rm old}+1)>$\\
\midrule
2&1&1&0.1&1&145&4.64\\
3&1&1&0.01&1&2,230&43.73\\
3&1&1&0.01&1/2&1,242,418&43.73\\
5&2&1&0.01&1&236,784&87.15\\
3&3&2&0.001&1&$6.949\cdot10^{14}$&2606.24\\
\bottomrule
\end{tabular}
\caption{Exact new upper exponents from \eqref{eq:Hcondition}, with
$h_0=2$ except $h_0=3$ in the last row. The old column is a rounded-down
lower bound on the exponent needed to satisfy the harmonic condition for
$B=nB_*$; it is not a lower bound on the true optimal order. The
$\eta=1/2$ comparison is with old existence, not old density.}
\end{table}

\begin{proposition}[A necessary budget dependence]\label{prop:lower}
For any $h_0\le L$, there is an integral polynomial with
$\Gauss_\rho=L-h_0+1$ (if $\rho\ge r_{h_0}$) that vanishes at all
primitive torsion points in shells $h_0,\ldots,L$.
\end{proposition}
\begin{proof}
Take $g(T)=\prod_{h=h_0}^{L}\Phi_{p^h}(1+T)$. Each factor has
$e_h$ roots all of valuation $r_h$, and is monic. Its Gauss valuation
at $\rho$ is $e_hr_h=1$. The vanishing is immediate. In fact
$v(g(0))=L-h_0+1$, so a finite witness of that valuation exists.
\end{proof}
This shows a linear dependence on the budget is unavoidable in some
regimes. We make no claim that the new polynomial exponents are sharp.
\section{Multivariate and analytic transfers}

\begin{theorem}[Finite-packet specialization]\label{thm:multi}
Let $F_i\in\Ocal[T_1,\ldots,T_s]$, $1\le i\le n$, and suppose that
for each $i$ there is a witness $x_i=(x_{i1},\ldots,x_{is})$ with
$v(x_{ij})\ge\rho>0$ and
\[
 t_i=v(F_i(x_i))\le B_i<\infty.
\]
Fix $s_0\ge s\ge1$ and tolerances $\epsilon_i>0$. For
Theorem \ref{thm:one}, use the uniform one-coordinate budgets
$B'_i=B_i+\epsilon_i$, $\delta_i=\epsilon_i/s_0$ and the same
$h_0,\eta$ at every step. Let $H$ satisfy \eqref{eq:Hcondition} for
these parameters. Then there are at least
\[
 K^s\ \text{distinct tuples},\qquad
 K=\lfloor(1-\eta)e_{h_0}\rfloor+1,
\]
of primitive roots $\zeta_j$ of orders $p^{h_j}$ with $h_0\le h_j\le H$
such that
\[
 v(F_i(\zeta_1-1,\ldots,\zeta_s-1))\le t_i+\epsilon_i
 \quad\text{for all $i$.}
\]
\end{theorem}
\begin{proof}
Replace coordinates one at a time. At any common prefix of replacements,
freeze all coordinates except the next in each polynomial at its own
current witness (the prefix is common). This gives integral univariate
polynomials $g_i$ whose values at the next witness coordinates have
valuation at most $t_i+(j-1)\delta_i$, so $g_i\ne0$ and
$\Gauss_\rho(g_i)\le B_i+\epsilon_i$. Theorem \ref{thm:one} finds
at least $K$ distinct common choices for the next coordinate, each
increasing the current witness valuation by at most $\delta_i$.
Keep those branches. After $s$ steps each has loss at most
$s\delta_i\le\epsilon_i$, and distinct prefixes produce distinct
complete tuples. Every chosen coordinate is nontrivial and lies in
the claimed range. The bound was chosen before seeing any polynomial.
\end{proof}

In particular, to obtain total positive excess at most $\epsilon$ for a
fixed packet one can choose $\epsilon_i=\epsilon/n$. The earlier
ProveIt report obtains this with an exponential exponent bound that is
independent of packet cardinality and even allows an infinite family.
Our improvement has a different trade-off: the polynomial estimate has
an exponent depending on $n$. Neither result subsumes the other.

\begin{corollary}[Bounded power series]\label{cor:series}
Theorem \ref{thm:multi} remains true for
$F_i\in\Ocal_{\C_p}[[T_1,\ldots,T_s]]$ and witnesses in $\C_p^s$,
with the same $H$ and count $K^s$.
\end{corollary}
\begin{proof}
The candidate torsion grid is finite; every coordinate in the grid has
valuation at least $r_H>0$. Choose $M>\max(\rho,\max_i(B_i+\epsilon_i))$.
Truncate the power series at total degree $d$ where $(d+1)r_H>M$.
The omitted tails have valuation $>M$ on the grid and at every witness.
Approximate the finitely many integral coefficients and witness
coordinates by algebraic integral numbers to valuation precision $>M$.
The witnesses retain coordinate valuations at least $\rho$. Evaluating
integral polynomials on integral points is $1$-Lipschitz for the maximum
$p$-adic norm (telescope monomials), so these changes have valuation
$>M$ in the values. The witness valuations are unchanged, and the
polynomial theorem gives $K^s$ tuples with selected valuations $<M$.
The original series have exactly those valuations at each selected
tuple by the strong triangle inequality. This proves the assertion.
\end{proof}

\begin{corollary}[Presentation matrices]\label{cor:matrix}
For finitely many matrices of bounded integral power series with full-row-rank
minors at witnesses, the same grid yields at least $K^s$ simultaneous
specializations where the chosen minors remain nonzero and satisfy the
valuation bounds of Theorem \ref{thm:multi}. If their coefficients lie
in a fixed finite extension $L/\Q_p$, the normalized lengths of the
corresponding cokernels over $L(\mu_{p^H})$ are bounded above by the
valuation of the chosen minors.
\end{corollary}
\begin{proof}
Apply Corollary \ref{cor:series} to the selected determinants. For a
square matrix of full rank over a discrete valuation ring, Smith normal
form identifies normalized cokernel length with determinant valuation.
The full rectangular presentation has cokernel a quotient of that of
the selected square submatrix, so its length is no larger. Here length
is normalized by dividing by the absolute ramification index.
\end{proof}

These statements concern bounded coefficients. The unbounded series
$\log(1+T)$ vanishes at every $p$-power torsion point but not at $p^2$;
there is no such selection rule for arbitrary analytic functions. Compare
the rigidity questions in \cite{Serban18,Serban22}, which concern torsion
zero sets and do not imply the finite budget estimate above by their
stated conclusions.

\section{Simultaneous powers and the arithmetic theorem}
\subsection{Controlling powers of one twist}
Suppose the $F_i$ and witnesses satisfy the hypotheses of Theorem
\ref{thm:multi}, with witness depth $r>0$. Let $m\ge1$. Apply that
theorem to the $nm$ polynomials
\[
 F_{i,j}(T)=F_i\big((1+T_1)^j-1,\ldots,(1+T_s)^j-1\big),
 \qquad 1\le j\le m.
\]
For each coordinate $x$ of a witness and integer $j\ge1$, there is an
algebraic $y$ satisfying $(1+y)^j-1=x$ and $v(y)\ge r/j$. To see this,
if $x=0$ take $y=0$; otherwise the polynomial
$(1+Y)^j-1-x$ is monic integral, so its roots are integral, and their
valuations sum to $v(x)\ge r$. At least one is at least $r/j$.
The new witnesses thus have depth $\rho=r/m$ and their values have
exactly the original valuations $t_i$. We take the same budgets
$B_i$ and tolerances $\epsilon_i$ for all $j$ copies of $i$.

Choose $h_0$ with $r_{h_0}\le r/m$ and $p^{h_0}>m$.
For every selected component $\zeta$ and every $1\le j\le m$,
\[
 v(\zeta^j-1)=p^{\vp(j)}v(\zeta-1)\le m r_{h_0}\le r,
\]
and no $\zeta^j$ equals one. The resulting $K^s$ tuples satisfy
$v(F_i((\zeta_u^j-1)_u))\le t_i+\epsilon_i$ for every $i,j$.
The bound $H$ is polynomial in inverse error tolerance for fixed $n,m,s_0$;
its exponent increases with $m$. This is the quantifier order used by
\texttt{exists\_rootsOfUnity\_ellVal\_eval\_le} in the source.

\subsection{Arithmetic strengthening under the source hypotheses}
We now use the following input from the pinned readable edition (see
\path{Line.lean}, lines 630--687, and \path{GeneralTheorem.lean},
lines 392--457). We state it in words to avoid misidentifying the
source's precise local hypotheses.

Fix the CM field $E$, prime $p$, set $S_0$ including the places above
$p$, inertia element $g_0$ at a place of $S_0$ away from $p$, coefficient
ring, Selmer data $T_i$, and Stickelberger systems $Y_k$ of the source.
Assume \texttt{IsTorsionLevel} at its required torsion level $M$ and all
fields of the \texttt{SelmerDatum} and \texttt{StickelbergerSystem}
structures. Write $\beta_a$ for their finite normalized valuation
budgets, including the global and local terms in the source. In
particular the source Selmer datum has special scalar inertia,
local cohomological conditions, and representation hypotheses; these
are not automatic for arbitrary elliptic curves.

The source constructs an integer $s_0\ge1$, depending only on the fixed
data, such that for all rational $r>0$ with $s_0r<1$ and all sufficiently
large $N$ there exist the following data:
\begin{itemize}
\item An auxiliary set $Q$ of degree-one primes, disjoint from its
conjugate $\overline Q$ and the prescribed bad sets, with
$2|Q|\le s_0$. In the proof $Q=J$ contains the designated first prime
$P_0$, so $Q$ can be chosen nonempty.
\item Characters $\lambda_u:G_E\to\Z/p^N\Z$ for $u\in Q$, unramified
outside $u,\bar u$ and trivial above $p$, whose restrictions to inertia
at $u$ are onto. At inertia at $u$, every other $\lambda_v$ is zero.
\item One integral polynomial $F_a$ per datum or system and a witness
of depth $r$ at which its valuation is at most $\beta_a+s_0r$.
\end{itemize}
For $p^N$-th roots of unity $(\zeta_u)_u$, the character
$\psi=\prod_u\zeta_u^{\lambda_u}$ is anticyclotomic. Polynomial
values at $(\zeta_u-1)_u$ evaluate the relevant group-ring elements or
bound the normalized twisted Selmer length, the latter when the
valuation is strictly less than $\beta_a+1$. This description is exactly
the source's \texttt{TwistPolynomials} construction. The finishing
lemma \texttt{isSmallTwist\_twist} gives the local and global conclusions
from the polynomial valuations and the coordinate bound.

\begin{theorem}[Bounded-order packets on one support]\label{thm:arith}
Under the source hypotheses just specified, fix $m\ge1$, $\epsilon>0$,
and a desired number $A\ge1$ of distinct characters. There is an
explicit upper exponent $H$ in terms of $p,m,\epsilon,s_0$ and the
finite budgets $\beta_a$, and $A$, with the following property.
One can choose a single nonempty auxiliary set $Q$ as above and at least
$A$ distinct anticyclotomic characters $\psi$ such that:
\begin{enumerate}
\item each $\psi$ has order dividing $p^H$;
\item every power $\psi^j$, $1\le j\le m$, is ramified at every
prime of $Q\cup\overline Q$ and unramified outside it, trivial at the
places above $p$, and has the same tame conductor
$\prod_{u\in Q\cup\overline Q}u$;
\item every one of those powers satisfies the source's
\texttt{IsSmallTwist} inequalities with tolerance $\epsilon$: normalized
Selmer lengths $\le\beta_i+\epsilon$, Stickelberger valuations
$\le\beta_k+\epsilon$, and the sum over the tame-generator valuations
$\le\epsilon$.
\end{enumerate}
For fixed arithmetic data and fixed $m,A$, the exponent $H$ can be
chosen polynomially bounded in $1/\epsilon$ as $\epsilon\downarrow0$.
There is no asserted bound here on the norms of the primes in $Q$.
\end{theorem}
\begin{proof}
Choose a rational $0<e\le\min(\epsilon,1/2)$ with
$e\ge\tfrac12\min(\epsilon,1/2)$, and rational $r=e/(4s_0)$.
Use the powered-polynomial construction with $B_a=\beta_a+1$,
$\epsilon_a=e/2$, $\rho=r/m$, $\eta=1/2$, and $s_0$ as the coordinate
bound. For all $m$ copies use the same budgets. Choose $h_0$ large enough
that $r_{h_0}\le r/m$, $p^{h_0}>m$, and
$K=\lfloor e_{h_0}/2\rfloor+1\ge A$. Choose $H$ by
\eqref{eq:Hcondition} or \eqref{eq:H} with the increased intermediate
budgets $B_a+e/2$ and losses $e/(2s_0)$. All these choices precede the
actual polynomials. If the packet has no constraints, the same selection
argument works trivially using arbitrary roots (or append $F=1$).

Next choose $N\ge H$ sufficiently large that the source construction
exists, with its designated first prime. This is possible by the
``all sufficiently large $N$'' conclusion. The polynomials $F_a$ have
witnesses with valuations $t_a\le\beta_a+s_0r<\beta_a+1$.
The multivariate theorem for powers now gives at least
$K^{|Q|}\ge A$ tuples whose every power has coordinate valuations
$\le r$ and polynomial valuations
\[
 v(F_a((\zeta_u^j-1)_u))\le t_a+e/2
 \le\beta_a+3e/4\le\beta_a+e<\beta_a+1.
\]
The source's finishing lemma applies with tolerance $e$, since
$s_0r=e/4\le e<1$; monotonicity increases the tolerance to $\epsilon$.
Every $\zeta_u$ is a $p^H$-th root and hence a $p^N$-th root. Thus
$\psi^{p^H}=1$, even if the internal auxiliary level $N$ is much larger
than $H$.

For injectivity, for each $u\in Q$ fix $g_u$ in inertia above $u$ with
$\lambda_u(g_u)=1$ and $\lambda_v(g_u)=0$ for $v\ne u$. Such an
element is chosen independently of the tuple. Then $\psi(g_u)=\zeta_u$,
so different tuples give different characters. In fact their orders
are exactly the maximum component order. Since each power has nontrivial
inertia at each $u$ and anticyclotomic conjugate inertia has the inverse
value, each power is also ramified at $\bar u$. Outside these primes
it is unramified; all these primes are away from $p$ and the characters
have $p$-power order, hence are tame. A nontrivial tame one-dimensional
character has conductor exponent one. This proves the common conductor
assertion. The polynomial asymptotic follows from the fixed number of
polynomials and the fact that $h_0$ grows only logarithmically with
$1/e$ for fixed $A,m$.
\end{proof}

\paragraph{Interpretation.}
The new count is for a common support (and hence a common conductor
ideal) chosen after specifying the packet size. It is not a positive
density statement as conductors are ordered by norm. The proof uses
one common auxiliary $N$, one common set $Q$, and one fixed character
parametrization for all successful tuples. The result adds an explicit
order cap and branching multiplicity to the source's qualitative
existence theorem. The existence of many characters alone can also be
approached through Galois orbits in particular arithmetic settings;
no novelty claim is made for such qualitative multiplicity in isolation.

\section{Implications, limits, and further research}

\paragraph{Central-value nonvanishing.}
Kriz--Nordentoft's horizontal $p$-adic $L$-function work \cite{KN}
relates integral measures and twisted central values and proves
simultaneous nonvanishing results. If a finite packet of central values
in a horizontal family is represented by bounded integral interpolation
functions times nonzero factors, and each function has a witness of
known depth and bounded valuation, Corollary \ref{cor:series} supplies
simultaneous nonvanishing and a finite grid with explicit order cap.
The theorem even yields a good fraction on a one-variable shell.
This is a concrete local reduction for quantitative questions of
higher-order twisted nonvanishing. It does not prove the
David--Fearnley--Kisilevsky prediction (or Goldfeld's conjecture):
fixed-order twists, density in conductor, functional-equation signs,
and arithmetic interpolation must be handled separately. In particular,
the selected order can vary within the grid. Proposition \ref{prop:lower}
explains why prescribing it in advance for arbitrary data is impossible.

\paragraph{Selmer and Iwasawa questions.}
Corollary \ref{cor:matrix} can control multiple presentation minors and
normalized lengths at the same finite-order specializations. To apply
it to a particular Iwasawa main conjecture one still needs the relevant
arithmetic presentation, control of local conditions under specialization,
and a nonzero witness; determinant identities are not a formal
consequence of the selector. The source's scalar-inertia restriction
also blocks direct application to the ordinary rank-two elliptic Tate
module: its determinant is cyclotomic, hence trivial on inertia away
from $p$; a scalar $\xi I_2$ there would satisfy $\xi^2=1$, contrary to
requiring $\xi^2-1$ to be a unit. A preliminary twist or a different
representation would require a new comparison argument.

\paragraph{Normalized length is not vanishing.}
If a coefficient field has absolute ramification index $e_L$, a
nonzero finite-length module has normalized length at least $1/e_L$,
not at least one. In a nonempty ramification packet with all powers up
to $m$, let $k=\lfloor\log_p m\rfloor$. The local sum constraint forces
$\epsilon e_L\ge2|Q|p^k$ by evaluating at the allowed power $p^k$ and
using $v(\zeta^{p^k}-1)\ge p^k/e_L$ on each nontrivial inertia value.
This is the prior report's ramification obstruction. Improving $H$
does not remove it. To deduce vanishing of a fixed ordinary Selmer
module, a descent comparison to the twisted modules with normalized
length preserved would be needed.

\paragraph{Research questions.}
(i) Can the exponent $n/(c\eta)$ be reduced, perhaps by accounting for
the exact saving $1-p^{-(h-k)}$ rather than its uniform lower bound $c$?
(ii) Is there a degree-free bound with polynomial dependence on the
total budget and tolerance, independent of the packet size, while
retaining total positive excess? (iii) Can one convert the shell density
into conductor-norm density by bounding the norms of the auxiliary primes
and tracking how the character parameter space varies? (iv) For a
specific horizontal interpolation system, can the local bound combine
with a witness and uniform Euler-factor estimates to improve a published
arithmetic nonvanishing theorem? (v) Can the small-twist argument be
reworked to admit rank-two representations without the special scalar
inertia? (vi) The source's compact height comparison is another possible
transfer route for nonabelian $L$-functions, but it needs analogues of its
short, spanning and long form inputs and is not resolved by the present
method.

\section{Reproducibility and formalization plan}

\texttt{verify.py} uses Python's standard library and exact rational
arithmetic. It checks the cyclotomic derivative identity and capped
averaging in small shells, tests actual products of linear factors at
primitive roots (including full-shell vanishing and random non-Galois
invariant multisets), checks each bad-shell potential inequality, and
calculates the least integer $H$ satisfying the displayed sufficient
condition for rational parameters. Valuations of pairwise differences
are exact: $v(\zeta_a-\zeta_b)=r_{h-v_p(a-b)}$ for distinct primitive
roots of the same order. No floating-point values enter the theorem
checks or the computation of $H$; floating point is used only to display
a logarithmic comparison in the table. \texttt{verification.txt}
contains the actual recorded run. The tests are not a proof of the
universal result; the mathematical proofs above supply that argument.

A Lean port can reuse \texttt{rootRadius} and \texttt{rootsOfValuation},\\
plus the identities \texttt{gaussVal\_eq\_add\_sum\_roots} and\\
\texttt{ellVal\_eval\_add\_card\_mul}. The new lemmas are (1) capped
excess averaging, (2) the slack lower bound from preceding shells,
(3) the finite bad-fraction implication, and (4) the product potential
contradiction. To avoid real powers in the core finite theorem, a
sufficient assumption is directly
$\prod_{h=h_0}^H(1+c\eta/(h+C))\ge P$, then prove the closed form by a
separate concavity lemma. Use the existing coordinate induction and
power-substitution arguments after replacing the univariate selection
step; preserve the fact that $H$ is chosen before the polynomials.
The arithmetic wrapper also needs to expose the designated first prime
and use \texttt{exists\_eq\_ofAdd\_one} to prove injectivity. See
\texttt{LEAN\_PORT.md} for a more explicit dependency map.
No new Lean file has been compiled; Lean and Lake are absent from this
execution environment. The source audit is based on retrieved GitHub
files, not a local checkout. The report therefore distinguishes the
human mathematical extension from a machine-checked extension.

\begin{thebibliography}{9}
\bibitem{upstream} L. Sernau and collaborators, \emph{GRH for Dirichlet
L-functions, in Lean 4}, repository based on the proof of L. Alp\"oge
with Claude (2026). \url{https://github.com/lsernau/grh-lean}.
Pinned commit \texttt{72fead102f663d5c7a407964a13dc0cf9d977207}.
\bibitem{KN} D. Kriz and A. C. Nordentoft, \emph{Horizontal $p$-adic
$L$-functions}, arXiv:2310.20678 (2023, revised 2024).
\url{https://arxiv.org/abs/2310.20678}.
\bibitem{Serban18} V. Serban, \emph{An infinitesimal $p$-adic
multiplicative Manin--Mumford conjecture}, J. Th\'eor. Nombres Bordeaux
30 (2018), 393--408. \url{https://doi.org/10.5802/jtnb.1030}.
\bibitem{Serban22} V. Serban, \emph{On $p$-adic versions of the
Manin--Mumford conjecture}, IMRN (2022), 10279--10301.
\url{https://doi.org/10.1093/imrn/rnab136}.
\bibitem{prior} ProveIt, incoming reports,
\url{https://github.com/VladimirReshetnikov/ProveIt/tree/main/docs/incoming}.
In particular the two archives identified in Section 1; they are
research drafts rather than independently verified published work.
\end{thebibliography}
\end{document}
